% ============================================================
\section{Espacios simplemente conexos — Teorema 4.4}
% ============================================================

\begin{frame}{Definición: Espacio simplemente conexo}

  \begin{block}{Definición}
    Un espacio \emph{arco-conexo} $X$ es \textbf{simplemente conexo}
    si $\pione(X)$ es el grupo trivial.
  \end{block}

  \bigskip
  \begin{itemize}
    \item Recordemos que $\pione(X, \xo)$ es el \textbf{grupo fundamental}
          de $X$ en el punto base $\xo$.
    \item Cuando $X$ es arco-conexo, el grupo abstracto resultante no
          depende de la elección de $\xo$, por lo que escribimos
          simplemente $\pione(X)$.
    \item Ser simplemente conexo significa que \emph{todo lazo} en $X$
          basado en cualquier punto es homotópico al lazo constante.
  \end{itemize}

\end{frame}

\begin{frame}{Recordatorio: Espacio contráctil}

  \begin{alertblock}{Definición (recordatorio)}
    Un espacio topológico $X$ es \textbf{contráctil} si la aplicación
    identidad $\mathrm{id}_X : X \to X$ es homotópica a una aplicación
    constante; es decir, existe un punto $\xo \in X$ y una homotopía
    \[
      H : X \times I \;\longrightarrow\; X
    \]
    tal que
    \[
      H(x,\,0) = x \qquad \text{y} \qquad H(x,\,1) = \xo
      \qquad \forall\, x \in X.
    \]
  \end{alertblock}

  \bigskip
  \textbf{Ejemplos:} $\R^n$, cualquier conjunto convexo, el cono sobre
  un espacio compacto, etc.

\end{frame}

\begin{frame}{Teorema 4.4}

  \begin{block}{Teorema 4.4}
    Todo espacio contráctil es simplemente conexo.
  \end{block}

  \bigskip
  \textbf{Estrategia de la demostración:}
  \begin{enumerate}
    \item Mostrar que $X$ es \textbf{arco-conexo}.
    \item Construir, para cada lazo $\alpha$ en $X$ basado en $\xo$,
          una homotopía $h : I \times I \to X$ que demuestre
          $\clase{\alpha} = \clase{c}$, donde $c$ es el lazo constante.
  \end{enumerate}

\end{frame}

\begin{frame}{Demostración — Parte 1: Arco-conexidad}

  Sea $X$ un espacio contráctil. Existe $\xo \in X$ y una homotopía
  \[
    H : X \times I \;\longrightarrow\; X
    \quad \text{con} \quad
    H(x,0) = x, \quad H(x,1) = \xo, \quad \forall\, x \in X.
  \]

  \medskip
  \begin{block}{$X$ es arco-conexo}
    Para cada $x \in X$, la función
    \[
      \alpha_x = H(x,\,\cdot\,) : I \;\longrightarrow\; X
    \]
    es un camino de $H(x,0) = x$ hasta $H(x,1) = \xo$.

    \medskip
    Dados $x, y \in X$ cualesquiera, el camino
    $\alpha_x * \bar{\alpha}_y$ une $x$ con $y$.
    Luego $X$ es arco-conexo. \hfill $\checkmark$
  \end{block}

\end{frame}

\begin{frame}{Demostración — Parte 2: Caso $H(\xo, s) = \xo$}

  Supongamos, por ahora, que $H$ satisface la propiedad adicional
  \[
    H(\xo,\, s) = \xo \qquad \forall\, s \in I.
  \]

  \medskip
  Sea $\clase{\alpha} \in \pione(X, \xo)$.
  Definimos la homotopía $h : I \times I \to X$ por
  \[
    h(t,\, s) \;=\; H\!\bigl(\alpha(t),\, s\bigr).
  \]

  \begin{exampleblock}{Verificación de $h$}
    \begin{align*}
      h(t,0) &= H(\alpha(t), 0) = \alpha(t), \\[4pt]
      h(t,1) &= H(\alpha(t), 1) = \xo, \quad t \in I, \\[4pt]
      h(0,s) &= H(\alpha(0),\, s) = H(\xo,\, s) = \xo, \\[4pt]
      h(1,s) &= H(\alpha(1),\, s) = H(\xo,\, s) = \xo, \quad s \in I.
    \end{align*}
  \end{exampleblock}

\end{frame}

\begin{frame}{Demostración — Parte 2: Conclusión}

  La homotopía $h$ satisface:
  \begin{itemize}
    \item $h(\cdot, 0) = \alpha$ \quad (nivel inicial es el lazo $\alpha$),
    \item $h(\cdot, 1) \equiv \xo$ \quad (nivel final es el lazo constante $c$),
    \item $h(0, s) = h(1, s) = \xo$ para todo $s \in I$
          \quad (los extremos permanecen fijos en $\xo$).
  \end{itemize}

  \bigskip
  \begin{block}{Conclusión}
    $h$ es una homotopía de lazos de $\alpha$ al lazo constante $c$. Por
    tanto $\clase{\alpha} = \clase{c}$, y $\pione(X, \xo)$ contiene
    únicamente el elemento identidad.

    \medskip
    Luego $X$ es \textbf{simplemente conexo}.
  \end{block}

\end{frame}

\begin{frame}{Demostración — Parte 3: Caso general}

  Si $H(\xo,\,\cdot\,) : I \to X$ \textbf{no} es el camino constante,
  la homotopía $h(t,s) = H(\alpha(t),s)$ ya no tiene sus extremos fijos
  en $\xo$.

  \medskip
  \textbf{Solución:} modificar cada nivel $s$ de $h$ para obtener
  un lazo basado en $\xo$.

  \medskip
  Sea $\beta_s = H(\xo,\,\cdot\,)|_{[0,s]}$ el camino recorrido por
  $\xo$ hasta el nivel $s$.
  La homotopía corregida concatena $\beta_s$, el nivel $s$ de $h$, y
  $\bar{\beta}_s$:
  \[
    \tilde{h}(\cdot,\, s) \;=\; \beta_s * h(\cdot,\, s) * \bar{\beta}_s.
  \]
  Esto produce un lazo basado en $\xo$ para cada $s$, con
  $\tilde{h}(\cdot,0) = \alpha$ y $\tilde{h}(\cdot,1) = c$.

  \medskip
  La construcción explícita se ilustra en la Figura~4.5 del texto. \hfill $\square$

\end{frame}

\begin{frame}{Figura 4.5 — Ilustración de la homotopía corregida}

  \begin{center}
  \begin{tikzpicture}[scale=1.1]
    % Espacio X (forma amorfa)
    \draw[thick, fill=gray!15]
      plot[smooth cycle, tension=0.7]
      coordinates {(-2.5,0) (-2,1.5) (-0.5,2.2) (1.5,2) (2.5,0.8)
                   (2,-1) (0,-1.8) (-1.5,-1.5)};

    % Punto base x0
    \filldraw[blue] (0,-0.8) circle (2pt) node[below, blue]
      {\small $\xo$};

    % Camino H(x0, .) desde x0
    \draw[blue, thick, ->, >=Stealth]
      (0,-0.8) .. controls (-0.3,0.2) and (-0.8,0.8) .. (-0.6,1.4)
      node[left] {\small $H(\xo,\cdot)$};

    % Punto H(x0, s) en el nivel intermedio
    \filldraw[blue!70] (-0.6,1.4) circle (2pt)
      node[above left, blue!80] {\small $H(\xo,s)$};

    % Lazo h(·,s) en el nivel s
    \draw[coralOscuro, thick]
      (-0.6,1.4) .. controls (0.2,2.0) and (1.2,1.8) ..
      (0.8,1.0) .. controls (0.4,0.4) and (-0.1,0.8) .. (-0.6,1.4);
    \node[coralOscuro] at (0.6,1.7) {\small $h(\cdot,s)$};

    % Indicación del nivel
    \node at (2.0,1.8) {\small $X$};
  \end{tikzpicture}
  \end{center}

  \vspace{-2mm}
  \begin{itemize}
    \item El camino $H(\xo,\cdot)$ (azul) puede no ser constante.
    \item El lazo corregido $\tilde{h}(\cdot,s)$ recorre $\beta_s$,
          el nivel $h(\cdot,s)$ (rojo), y regresa por $\bar{\beta}_s$.
  \end{itemize}

\end{frame}

\begin{frame}{Resumen — Teorema 4.4}

  \begin{block}{Teorema 4.4 (reiteración)}
    Todo espacio contráctil es simplemente conexo.
  \end{block}

  \bigskip
  \textbf{Pasos clave de la demostración:}
  \begin{enumerate}
    \item La contracción $H$ proporciona caminos $\alpha_x = H(x,\cdot)$
          que unen cada punto con $\xo$, probando la \textbf{arco-conexidad}.
    \item Se define $h(t,s) = H(\alpha(t),s)$ para mostrar que cualquier
          lazo $\alpha$ en $\xo$ es homotópico al lazo constante.
    \item Si $H(\xo,\cdot)$ no es constante, se \textbf{corrige} $h$
          concatenando el camino $H(\xo,\cdot)$ en cada nivel.
    \item En ambos casos $\pione(X,\xo) = \{[c]\}$: el grupo trivial.
  \end{enumerate}

  \bigskip
  \begin{exampleblock}{Corolario inmediato}
    $\R^n$, los conjuntos convexos y todos los contráctiles tienen
    grupo fundamental trivial.
  \end{exampleblock}

\end{frame}
